\section{Maintaining Constraints Automatically}
\label{sec:maintain_constraints}

This section explains how constraints in Babelsberg/S are maintained automatically. 

An application in Babelsberg/S reacts to a write access on a constraint variable by solving the constraints automatically. Internally, a write access on a constraint variable calls \texttt{suggestValue:} on the variable. This will add a temporary equality constraint for the new value to the underlying constraint solver. The solver tries to solve all constraints. If the constraints are satisfiable, the value of all involved variables changes without notice to the application. Otherwise, an error is thrown and the new value is ignored. 

For example, if a Sudoku cell changes its value, all Sudoku constraints are solved in background. In other words, each change of a Sudoku cell adjust the whole Sudoku, depending on the constraints. That is possible, because the Sudoku application interacts with the underlying constraints.

Babelsberg/S handles different constraint solver. Currently, it supports Cassowary\cite{badros-tochi-2001}, which is implemented in Squeak, and Z3\cite{de2008z3}, which is an external program.